% Source: https://artofproblemsolving.com/community/c3h2704451_six_problems
% Stable identity: aops-topic-2704451-q02
By the midsegment theorem in triangles $ABD$, $BCD$, $ABC$, $ACD$:
\[PQ\parallel AB,\ PQ=6;\qquad QR\parallel CD,\ QR=5;\qquad RS\parallel AB,\ RS=6;\qquad SP\parallel CD,\ SP=5.\]
So $PQRS$ is a parallelogram with sides $6$ and $5$ parallel to $AB$ and $CD$. Since $\angle ABC+\angle BCD=150^\circ$, lines $AB$ and $CD$ meet at an angle of $180^\circ-150^\circ=30^\circ$. Hence
\[[PQRS]=6\cdot5\cdot\sin30^\circ=\boxed{15},\]
whatever the length of $BC$.

\begin{asy}
usepackage("amsmath");
size(4cm);
defaultpen(fontsize(10pt));
pair B=(0,0),C=(10,0),A=12*dir(80),D=C+10*dir(110);
pair P=(A+D)/2,Q=(B+D)/2,R=(B+C)/2,Sm=(A+C)/2;
draw(A--B--C--D--cycle,linewidth(0.8));draw(A--C,gray);draw(B--D,gray);
fill(P--Q--R--Sm--cycle,lightgray);draw(P--Q--R--Sm--cycle,linewidth(0.8));
label("$A$",A,N);label("$B$",B,SW);label("$C$",C,SE);label("$D$",D,NE);
label("$P$",P,N,fontsize(8pt));label("$Q$",Q,W,fontsize(8pt));label("$R$",R,S,fontsize(8pt));label("$S$",Sm,E,fontsize(8pt));
label("$80^\circ$",B+(0.9,0.6),fontsize(7pt));label("$70^\circ$",C+(-1.1,0.6),fontsize(7pt));
\end{asy}
